% running-headers — body text across two pages under a running header and a folio.
%
% Every page repeats the same header line and a page number in the footer. Neither is
% part of the document's text, so the gold file holds the body alone. The paragraph
% the page break cuts through must come out whole: removing the furniture is only
% half the job if the text either side of it stays split. Club and widow penalties
% are off so TeX puts the break inside the third paragraph, not before it. The page has
% ordinary margins, its header and folio in its top and bottom eighth as a page's are, and
% the body is set large so a short text still runs over the break.
\documentclass{article}
\usepackage{lmodern}
\usepackage{fancyhdr}
\usepackage[paperwidth=5in,paperheight=6in,top=0.9in,bottom=0.9in,hmargin=0.6in,headsep=0.3in,footskip=0.45in]{geometry}
\pagestyle{fancy}
\fancyhf{}
\fancyhead[C]{Field Notes on the Northern Lakes}
\fancyfoot[C]{\thepage}
\renewcommand{\headrulewidth}{0pt}
% No hyphenation: that is the hyphenation fixture's concern.
\hyphenpenalty=10000 \exhyphenpenalty=10000 \sloppy
\clubpenalty=0 \widowpenalty=0
\begin{document}
\fontsize{20}{24}\selectfont
The first lake is the largest, and the only one with a road along its shore. Its
water is clear to a depth of several metres in late summer.

The second lake is shallow and warm, and reeds cover its northern end.

The third lake lies in a hollow between two ridges, and the path to it climbs for an
hour through pine forest before it drops to the water. Few visitors make the walk, so
the shore is quiet even in August, and the only sound is the wind in the pines.

The fourth lake is the smallest, and it freezes first each winter.
\end{document}
